\documentclass[12pt]{article}
\usepackage[utf8]{inputenc}

\usepackage{setspace}
\usepackage[hypertexnames=false]{hyperref}
\usepackage[T1]{fontenc}

\usepackage{graphicx}
\usepackage{amsmath}
\usepackage{datetime}
\usepackage{comment}
\graphicspath{ {images/}} %upload your signature to this file
%Change the margins to fit your CV/resume content
\usepackage[left=1in, right=1in, bottom=.8in, top=.8in]{geometry}

%Changes the page numbers - {arabic}=arabic numerals, {gobble}=no page numbers, {roman}=Roman numerals
\pagenumbering{gobble}

\newdateformat{monthyeardate}{%
  \monthname[\THEMONTH], \THEYEAR}



%%%%%%%%%%%%%%%%% END OF PREAMBLE %%%%%%%%%%%%%%%%%%%%%
\begin{document}
\noindent

%%%%%%%%%%%%%%%%% NAME  %%%%%%%%%%%%%%%%%%%
\noindent\LARGE{\textsc{\textbf{Ruimin AO}}} \\
\vspace{0.2em}
\normalsize

\noindent Email: \href{mailto:R.ao@ucl.ac.uk}{R.ao@ucl.ac.uk}\\
% Phone: +44 07840207372 \\
\noindent \href{https://ruiminao.com/}{https://ruiminao.com/}
\vspace{1em}
\linespread{1.2}

%%%%%%%%%%%%%%%%% Education  %%%%%%%%%%%%%%%%%%%
\noindent\large{\textbf{\textsc{Education}}}
\vspace{0.2em}
\hrule
\normalsize
\noindent \begin{tabular}{ p{2cm} p{13.6cm}}
2020 -  & \textbf{PhD Economics}, University College London \\
2019 - 20 & \textbf{MRes Economics}, University College London \\
2018 - 19 & \textbf{MSc Economics with Distinction}, University of Warwick \\
2014 - 18 & \textbf{BSc Economics with Honours}, University of Nottingham \\
\end{tabular}
\vspace{0.6em}

%%%%%%%%%%%%%%%%% Affiliations  %%%%%%%%%%%%%%%%%%%
\noindent\large{\textbf{\textsc{Affiliations}}}
\vspace{0.2em}
\hrule
\normalsize
\noindent \begin{tabular}{ p{2cm} p{13.6cm}}
2023 -  & \textbf{Visiting PhD Student}, Institute for Fiscal Studies 
\end{tabular}
\vspace{0.6em}

%%%%%%%%%%%%%%%%% References  %%%%%%%%%%%%%%%%%%%
\noindent\large{\textbf{\textsc{Supervisors}}}
\vspace{0.2em}
\hrule
\normalsize
\noindent \begin{tabular}{ p{9cm} p{7cm}}
\textbf{Professor Marcos Vera-Hern\'andez}      & \textbf{Professor Val\'erie Lechene} \\
Department of Economics                         & Department of Economics     \\
University College London                       & University College London   \\
m.vera@ucl.ac.uk                                & v.lechene@ucl.ac.uk           
\end{tabular}
\vspace{0.6em}

%%%%%%%%%%%%%%%%% Field  %%%%%%%%%%%%%%%%%%%
\noindent\large{\textbf{\textsc{Field of Research}}}
\vspace{0.2em}
\hrule
\normalsize
\vspace{.8em}
\noindent{Health Economics, Applied Economics, Labour Economics, Behavioural Economics, Development Economics.}
\vspace{0.6em}


% %%%%%%%%%%%%%%%%% Work in Progress  %%%%%%%%%%%%%%%%%%%
% \noindent\large{\textbf{\textsc{Data Processing Experience}}}
% \vspace{0.2em}
% \hrule
% \normalsize
% \vspace{.8em}
% \noindent{I have rich experience and techniques in cleaning and handling large scale survey data. Throughout my PhD years I have used numerous large dataset from around the world including:} 
% 
% \textbf{Kantar Worldpanel}: the largest consumer behavior survey on representative households in the UK provides detailed product-level data on the spending habits of each household, with a substantial sample size of 30,000 households each year. Leveraging geographic information, I mapped households to different local economic shocks and analyzed how households responded to these variations. This analysis yields valuable insights for guiding effective and evidence-based policymaking.
% 
% \vspace{0.4em}
% \textbf{British Cohort Study, Next Step, Millennium Cohort Study}: These three cohort studies follow a representative sample of individuals born in certain years, where surveys were continuously carried out through their lifetime. I used these datasets to estimate a decision model for smoking behaviours, and used the results to conduct counterfactual studies to provide insights and evidence for policy making. 
% 
% \vspace{0.4em}
% \textbf{Consumer Pyramids Household Survey}: one of the largest longitudinal survey in India with around 160,000 household surveyed each wave. It provides detailed information on households and individuals, from labour force participation to health and well-being. I used this data to study the impact of the Covid-19 shock.
% 
% \vspace{0.4em}
% \textbf{UK Understanding Society}: the largest and most comprehensive household longitudinal survey in the UK, with more than 40,000 households surveyed in each wave. I utilised this data to examine how people's health changes after local and regional economic shocks.
% 
% \vspace{0.6em}

%%%%%%%%%%%%%%%%% Work in Progress  %%%%%%%%%%%%%%%%%%%
\noindent\large{\textbf{\textsc{Work in Progress}}}
\vspace{0.2em}
\hrule
\normalsize
\vspace{.8em}
\noindent$\bullet$~\textit{Behaviour Under Addiction: Dynamic Choice Between Smoking and Vaping.}\par
\vspace{.2em}
{\small\setlength{\leftskip}{1.5em}\setlength{\rightskip}{1.5em}
\noindent I develop and estimate a dynamic discrete choice model of joint cigarette and e-cigarette use, built around two distinct but interacting product-specific addiction stocks. The two products are substitutes for all consumers, yet the strength of substitution varies sharply with prior addiction: a smoker's sensitivity to cigarette prices weakens as smoking addiction deepens, whereas a vaper's sensitivity to e-cigarette prices \emph{strengthens} as vaping addiction deepens. The central finding is a striking asymmetry between the two taxes---a cigarette tax does little to curb smoking among the heaviest smokers, while an e-cigarette tax sharply reduces vaping but pushes former smokers and dual users back toward combustible cigarettes.\par}

\vspace{.4em}
\noindent$\bullet$~\textit{When climate change scrambles food prices, who sacrifices nutrition?}

\vspace{0.4em}
\noindent$\bullet$~\textit{The impact of discount retail stores on household consumption and wellbeing: evidence from the UK.}

\vspace{0.4em}
\noindent$\bullet$~\textit{Could doctor training improve knowledge and practice? An analysis of an RCT program in Yunnan, China.}

\vspace{0.6em}

% \end{comment}

\newpage


 %%%%%%%%%%%%%%%%% Grad Course  %%%%%%%%%%%%%%%%%%%
 \noindent\large{\textbf{\textsc{PhD Courses Taken}}}
 \vspace{0.2em}
 \hrule 
 \normalsize
 \noindent \begin{tabular}{ p{2cm} p{13.6cm}}
 ECON0107        & Macoreconomics \\
                & \textit{Course lecturer: Dr. Ralph Luetticke, Prof. Frank Portier, Prof. Vincent Sterk, Prof. Morten Ravn} \\
 ECON0106        & Microeconomics\\
                 & \textit{Course lecturer: Prof. Martin Cripps, Prof. Philippe Jehiel, Prof. Vasiliki Skreta} \\
 ECON0108        & Econometrics\\
                 & \textit{Course lecturers: Prof. Andrew Chesher, Prof. Toru Kitagawa, Prof. Martin Weidner and Dr. Daniel Wilhelm} \\
 ECON0086        & Topics in Applied Econometrics\\
                 & \textit{Course lecturers: Prof. Aureo de Paula} \\
 ECON0112        & Experimental Economics\\
                 & \textit{Course lecturers: Prof. Antonio Guarino, Dr. Terri Kneeland, Dr. Ryan Kendall} \\
 ECONDEV         & Development Economics\\
                 & \textit{Course lecturers: Prof. Imran Rasul, Prof. Jonas Hjort, Dr. Gabriel Ulyssea} \\
 ECON            & Topics in Applied consumer theory\\
                 & \textit{Course lecturer: Prof. Val\'erie Lechene} \\
 \end{tabular}
 \vspace{0.6em}
 
 % \begin{comment}

%%%%%%%%%%%%%%%%% Teach Exp  %%%%%%%%%%%%%%%%%%%
\noindent\large{\textbf{\textsc{Teaching Experience}}}
\vspace{0.2em}
\hrule 
\normalsize
\noindent \begin{tabular}{ p{2cm} p{13.6cm}}
2023 -          & TA for (BSc) Topics in Household Economics \\
                & \textit{Course lecturer: Prof. Val\'erie Lechene} \\
2021 -          & TA for (MSc) Health Economics \\
                & \textit{Course lecturer: Prof. Marcos Vera-Hern\'andez} \\
2020 -  23      & TA for (BSc) Applied Economics\\
                & \textit{Course lecturer: Dr. Dunli Li} \\
2020 - 21       & TA for (BSc) Quant. Economics \& Econometrics\\
                & \textit{Course lecturers: Prof. Dennis kristensen, Aureo de Paula and Kirill Borusyak} \\
\end{tabular}
\vspace{0.6em}

% \begin{comment}


%%%%%%%%%%%%%%%%% Work Exp  %%%%%%%%%%%%%%%%%%%
\noindent\large{\textbf{\textsc{Work Experience}}}
\vspace{0.2em}
\hrule 
\normalsize
\noindent \begin{tabular}{ p{2cm} p{13.6cm}}
2022 -          & Research Assistant to Prof. Marcos Vera-Hern\'andez \\
% 2017            & Development \& Planning Assistant, Ping An Health Cloud Ltd. \\
2015 - 16     & Research Assistant to Prof. Chang Liu\\
\end{tabular}
\vspace{0.6em}

%%%%%%%%%%%%%%%%% Awards  %%%%%%%%%%%%%%%%%%%
\noindent\large{\textbf{\textsc{Scholarships and Awards}}}
\vspace{0.2em}
\hrule 
\normalsize
\noindent \begin{tabular}{ p{2cm} p{13.6cm}}
2022            & Best TA on 1st Year UG module (UCL) \\
2018            & Outstanding Student Leader (UNNC) \\
2018            & Outstanding Graduate, by Government of Zhejiang Province, China \\
2017            & Head’s Scholarship (UNNC)\\
2017            & Guo-Hua Xing-Zhi Scholarship (UNNC) \\
2017            & Outstanding Youth League Leader (UNNC) \\
2016            & Dean’s Scholarship (UNNC) \\
2015            & Provost Scholarship (UNNC) \\
2014            & Entrance Scholarship (UNNC) \\
\end{tabular}
\vspace{0.6em}

%%%%%%%%%%%%%%%%% Others%%%%%%%%%%%%%%%%%%%
\noindent\large{\textbf{\textsc{Additional Information}}}
\vspace{0.2em}
\hrule
\normalsize
\noindent \begin{tabular}{ p{5cm} p{13.6cm}}
\textbf{Language}       & Mandarin (native), English (fluent) \\
\textbf{Programing Lanuages}      & Stata, \LaTeX, Matlab, R, Python\\
%\textbf{Citizenship}    & Chinese \\
\end{tabular}

% \vspace{1em}
% \noindent Updated: \monthyeardate\today
\end{document}
